In den letzten Jahren besteht die Möglichkeit für Forschungseinrichtungen und somit für die Lehre verwendete, kleine Satelliten zu bauen und in den Orbit zu schicken. Ermöglicht wurde dies durch die 1999 an der \acf{Cal Poly} entwickelten 	CubeSats.
\newline
Diese können für eine Vielzahl an Forschungsprojekten, jedoch auch für Telemedizinische Zwecke eingesetzt werden. 
So könnten beispielsweise Medizinische Daten wie zum Beispiel EKG-Signale mittels CubeSats von schwer erreichbaren Gegenden der Erde zu experten an weit entfernten Orten gesendet werden.
\newline
Da CubeSats sich im \ac{LEO} befinden, besteht teilweise nur ein 15 minütiges Zeitfenster in dem der Satellit das \ac{MCC} passiert, welches zu einer sehr geringen Datenausbeute führt \cite{preindl_optimization_2010}. Um mehr Daten zu erhalten, wurde \ac{GENSO} entwickelt, ein weltweites Netzwerk aus Bodenstationen welche nun nichtmehr nur auf deren eigene CubeSats sondern, mit Berechtigung der \acp{MCC}, auch auf andere zugreifen können \cite{leveque_global_2007}.
\newline
Diese Arbeit beschäftigt sich mit der Entwicklung einer neuen Open Source Alternative, \ac{BOINSO}, welche mithilfe von bereits vorhandenen Open Source Projekten und der Infrastruktur der \ac{BOINC} realisiert wird \cite{ries_boinc_2012}.
\newline
\ac{BOINSO} wird unter der \ac{GNU GPL} vertrieben und steht unter \url{https://github.com/Mehnen/BOINSO} zum Download bereit.